\documentclass[10pt,letterpaper]{article}

\usepackage[margin=0.72in]{geometry}
\usepackage[T1]{fontenc}
\usepackage{lmodern}
\usepackage{xcolor}
\usepackage[hidelinks]{hyperref}
\usepackage{microtype}
\usepackage{parskip}

\definecolor{accent}{HTML}{176B70}
\pagestyle{empty}
\setlength{\parindent}{0pt}
\setlength{\parskip}{7pt}

\begin{document}

{\Large\bfseries\color{accent} Piter Z. Garcia Bautista}\\[-1pt]
Rochester, NY \enspace|\enspace Remote\\
\href{mailto:garciapiterz@gmail.com}{garciapiterz@gmail.com}
\enspace|\enspace +1 (646) 288-3300\\
\href{https://linkedin.com/in/garciapiterz}{linkedin.com/in/garciapiterz}
\enspace|\enspace
\href{https://github.com/pzg8794}{github.com/pzg8794}

\vspace{8pt}
August 25, 2026

HMH / NWEA\\
Hiring Team

\textbf{Re: Senior Data Analyst, Requisition 20774}

Dear HMH and NWEA Hiring Team,

I am applying for the Senior Data Analyst position because it brings together
two parts of my work that I care about deeply: building reliable, reproducible
data workflows and helping educators make better-informed decisions for
students. I am completing M.S. degrees in Data Science at Rochester Institute
of Technology and Teaching Computer Science K--12 at the University of
Rochester in December 2026. I previously earned an M.S. in Computer Science and
bring more than ten years of experience across software, data systems, and
applied analytics.

At VEDADATA, I designed SQL, Python, and pandas workflows for data ingestion,
cleaning, transformation, validation, and analytics-ready delivery. I built
quality checks to identify missing, inconsistent, and unusual information and
worked with technical and nontechnical colleagues to clarify requirements,
document findings, and improve reliability. At VIOME, I developed
healthcare-oriented machine-learning workflows involving preprocessing,
feature engineering, model experimentation, evaluation, and AWS-based data
preparation. These roles taught me to treat data quality, provenance, and clear
documentation as core parts of analysis rather than afterthoughts.

My current RIT graduate assistantship adds a rigorous research dimension. I
build reproducible Python evaluation workflows for adaptive decision systems,
using structured logging, shared datasets, automated validation, and
comparison-ready reporting. I translate research questions and experimental
requirements into data structures, analyses, visualizations, and technical
artifacts that collaborators can inspect and reuse. This work closely matches
the position's emphasis on transparent pipelines, quality assurance,
troubleshooting, and partnership with researchers.

My University of Rochester teaching preparation gives me an additional lens on
the downstream use of education data. In K--5 computer science and community
science settings, I have used formative checks and student participation
evidence to adjust grouping, pacing, directions, and supports. That experience
reinforces why assessment and product data must be accurate, clearly defined,
and communicated with appropriate limitations when it informs decisions about
learners and instruction.

I would welcome the opportunity to support NWEA's Research and Policy
Partnerships team by making complex datasets more accessible, reliable, and
reusable. I would bring strong SQL and Python skills, a consistent focus on
data integrity and reproducibility, and the ability to communicate across
research, technical, and education contexts.

Thank you for considering my application. I would be glad to discuss how my
data-workflow, research, and teaching experience can contribute to HMH and
NWEA.

Sincerely,\\[10pt]
Piter Z. Garcia Bautista

\end{document}
